\documentclass[10pt]{article}
\usepackage[T1]{fontenc}
\usepackage[utf8]{inputenc}
\usepackage[french]{babel}
\usepackage{amsmath,amssymb,mathrsfs,mathpazo}
\usepackage{geometry}
\geometry{a4paper,margin=2cm}
\usepackage{graphicx,array,tabularx,tabularray,colortbl,multirow,multicol,booktabs,enumitem,pifont,ulem,fancybox,diagbox,cancel}
\usepackage{pstricks,pstricks-add,pst-plot,pst-tree,pst-node,pst-text,pst-func,pst-math,pst-bspline,pst-3dplot,pst-eucl,pst-solides3d}
\let\pscalculate\undefined
\usepackage{pst-bezier}
\usepackage{tikz,pgfplots,tkz-tab}
\usetikzlibrary{arrows,arrows.meta,patterns,calc,positioning,trees,decorations.markings,shapes.geometric,shapes.misc,shapes.arrows,intersections,decorations.pathmorphing,decorations.pathreplacing,backgrounds,shadows,babel}
\pgfplotsset{compat=1.18}
\usepgfplotslibrary{fillbetween}
\usepackage[np]{numprint}
\usepackage{listings,makecell,pas-tableur,twoopt}
\let\barre\undefined
\usepackage[pstricks]{tablvar}
\definecolor{darkgreen}{RGB}{0,100,0}
\providecommand{\red}{\color{red}}
\providecommand{\blue}{\color{blue}}
\providecommand{\green}{\color{green}}
\providecommand{\black}{\color{black}}
\providecommand{\white}{\color{white}}
\providecommand{\gray}{\color{gray}}
\providecommand{\cyan}{\color{cyan}}
\providecommand{\magenta}{\color{magenta}}
\providecommand{\yellow}{\color{yellow}}
\providecommand{\violet}{\color{violet}}
\providecommand{\orange}{\color{orange}}
\providecommand{\euro}{\texteuro}
\providecommand{\EUR}{\texteuro}
\providecommand{\ieme}{\textsuperscript{e}}
\providecommand{\degres}{\ensuremath{^\circ}}
\providecommand{\degree}{\ensuremath{^\circ}}
\providecommand{\vv}[1]{\overrightarrow{#1}}
\providecommand{\Coord}[1]{#1}
\providecommand{\decosix}{}
\providecommand{\starredbullet}{\ensuremath{\bullet}}
\pagestyle{empty}
\setlength{\parindent}{0pt}
\usepackage{adjustbox}\DeclareUnicodeCharacter{2236}{:}\DeclareUnicodeCharacter{2190}{\ensuremath{\leftarrow}}
\providecommand{\R}{\mathbb{R}}
\providecommand{\N}{\mathbb{N}}
\providecommand{\Z}{\mathbb{Z}}
\providecommand{\Q}{\mathbb{Q}}
\providecommand{\D}{\mathbb{D}}
\providecommand{\C}{\mathbb{C}}
\providecommand{\e}{{\mathrm{e}}}
\providecommand{\vect}[1]{\overrightarrow{#1}}
\providecommand{\vectt}[1]{\overrightarrow{\mathrm{#1}}}
\providecommand{\vvt}[1]{\overrightarrow{\mathrm{#1}}}
\providecommand{\np}[2][]{{#2~#1}}
\providecommand{\vv}[1]{{\overrightarrow{#1}}}
\renewcommand{\d}{{\mathrm{d}}}
\providecommand{\degres}{{^\circ}}
\providecommand{\degre}{{^\circ}}
\providecommand{\blue}{\color{blue}}
\providecommand{\red}{\color{red}}
\providecommand{\enskip}{\quad}
\providecommand{\llbracket}{\lbrack\!\lbrack}
\providecommand{\rrbracket}{\rbrack\!\rbrack}
\providecommand{\euro}{{\text{€}}}
\providecommand{\no}{{\text{n°}}}
\providecommand{\interval}[2]{{[#1;#2]}}
\providecommand{\intervFO}[2]{{[#1;#2[}}
\providecommand{\dotfill}{\ldots}

\providecommand{\ds}{\displaystyle}
\providecommand{\textding}[1]{#1}
\providecommand{\Oij}{\(\mathrm{O};\vec{\imath},\vec{\jmath}\)}
\providecommand{\Oijk}{\(\mathrm{O};\vec{\imath},\vec{\jmath},\vec{k}\)}
\providecommand{\Ouv}{\(\mathrm{O};\vec{u},\vec{v}\)}
\providecommand{\intervFF}[2]{\ensuremath{[#1;#2]}}
\renewcommand{\euro}{\ifmmode\text{\texteuro}\else\texteuro\fi}
\DeclareUnicodeCharacter{00B0}{\ifmmode^{\circ}\else\textdegree\fi}
\DeclareUnicodeCharacter{20AC}{\ifmmode\text{\texteuro}\else\texteuro\fi}
% Énoncés : Éducation nationale. Transcription et corrigés : APMEP, Denis Vergès et contributeurs. Découpage et classement : CoopMaths, Sébastien Lozano, Éric Elter et contributeurs. Adaptation de la présentation : Maths Explorer.
% CC BY-SA : https://coopmaths.fr/www/about/
% https://coopmaths.fr/alea/static/e3c/2020/tex/e3c_2020_00_sujet36_2_cor.tex
\begin{document}
\noindent\textit{APMEP / CoopMaths. CC BY-SA. Présentation adaptée par Maths Explorer.}\par\medskip
\subsection*{1.}

\paragraph{a.} Avec \(\overrightarrow{AB} = \begin{pmatrix} -2 \\ 4 \end{pmatrix}\) et \(\overrightarrow{AC} = \begin{pmatrix} 1 \\ 2 \end{pmatrix}\), on a :
\[
\overrightarrow{AB} \cdot \overrightarrow{AC} = -2 \times 1 + 4 \times 2 = -2 + 8 = 6.
\]

\paragraph{b.} On a :
\[
\|\overrightarrow{AB}\|^2 = \overrightarrow{AB} \cdot \overrightarrow{AB} = 4 + 16 = 20, \text{ d'où } \|\overrightarrow{AB}\| = \sqrt{20} = \sqrt{4 \times 5} = 2 \sqrt{5},
\]
\[
\|\overrightarrow{AC}\|^2 = \overrightarrow{AC} \cdot \overrightarrow{AC} = 1 + 4 = 5, \text{ d'où } \|\overrightarrow{AC}\| = \sqrt{5}.
\]

\paragraph{c.} On peut également écrire le produit scalaire :
\[
\overrightarrow{AB} \cdot \overrightarrow{AC} = \|\overrightarrow{AB}\| \times \|\overrightarrow{AC}\| \times \cos(\widehat{BAC}),
\]
soit :
\[
6 = 2 \sqrt{5} \times \sqrt{5} \times \cos(\widehat{BAC}) \iff \cos(\widehat{BAC}) = \dfrac{6}{2 \sqrt{5} \times \sqrt{5}} = \dfrac{6}{10} = \dfrac{3}{5} = 0{,}6.
\]
La calculatrice donne \(\widehat{BAC} \approx 53{,}1^\circ\), soit \(53^\circ\) au degré près.

\subsection*{2.}

\paragraph{a.} Une équation de la droite \((AB)\) étant \(y = ax + b\), alors on a le système :
\[
\begin{cases} 
-1 = 2a + b \\ 
3 = 0 \times a + b 
\end{cases}
\]
d'où \(b = 3\) et, par substitution, on obtient \(-1 = 2a + 3\) ou \(a = -2\).

Donc :
\begin{align*}
&M(x\,;\,y) \in (AB) \\
\iff &y = -2x + 3 \\
\iff &2x + y - 3 = 0.
\end{align*}

\paragraph{b.} Si \(H(x\,;\,y)\), on a \(\overrightarrow{CH} \begin{pmatrix} x - 3 \\ y - 1 \end{pmatrix}\), et les coordonnées vérifient :
\begin{align*}
&\begin{cases} 
H \in (AB) \\ 
\overrightarrow{CH} \cdot \overrightarrow{AB} = 0
\end{cases} \\
\iff &\begin{cases} 
2x + y - 3 = 0 \\ 
-2(x - 3) + 4(y - 1) = 0 
\end{cases} \\
\iff &\begin{cases}
2x + y - 3 = 0 \\
-2x + 6 + 4y - 4 = 0
\end{cases} \\
\iff &\begin{cases}
2x + y - 3 = 0 \\
-2x + 4y + 2 = 0
\end{cases}
\end{align*}

Ce qui donne par résolution :
\[
5y - 1 = 0 \iff y = \dfrac{1}{5},
\]
et :
\begin{align*}
&2x + \dfrac{1}{5} - 3 = 0 \\
\iff &2x - \dfrac{14}{5} = 0 \\
\iff &2x = \dfrac{14}{5} \\
\iff &x = \dfrac{7}{5}.
\end{align*}
Le point \(H\) a donc pour coordonnées \(H\left( \dfrac{7}{5}\,;\,\dfrac{1}{5} \right)\).

\end{document}
